\documentclass[11pt]{article}

\usepackage[margin=1in]{geometry}
\usepackage{amsmath,amssymb,amsthm}
\usepackage{mathtools}
\usepackage{empheq}
\usepackage{bm}
\usepackage{enumitem}
\usepackage{hyperref}

\newcommand{\E}{\mathbb{E}}
\newcommand{\R}{\mathbb{R}}
\newcommand{\X}{\mathcal{X}}
\newcommand{\Tr}{\mathcal{T}}
\newcommand{\Tstar}{\mathcal{T}^{*}}
\newcommand{\Gop}{\mathcal{G}}
\newcommand{\Aop}{\mathcal{A}}
\newcommand{\Sop}{\mathcal{S}}
\newcommand{\gss}{g^{\mathrm{ss}}}
\newcommand{\Vss}{V^{\mathrm{ss}}}
\newcommand{\css}{c^{\mathrm{ss}}}
\newcommand{\sss}{s^{\mathrm{ss}}}
\newcommand{\dd}{\,\mathrm{d}}

\title{The Krusell--Smith Economy in Discrete Time:\\ A Master Equation Formulation}
\author{}
\date{\today}

\begin{document}
\maketitle

\begin{abstract}
Bilal (2023, ``Solving Heterogeneous Agent Models with the Master Equation'')
develops the Master Equation representation of heterogeneous-agent economies and
its first-order approximation (the FAME) in \emph{continuous} time, using the
Krusell and Smith (1998) economy as the leading example. This note writes down
the \emph{discrete-time} analogue. The logic is unchanged: the value function is
defined on the space of distributions and a single Bellman equation summarizes
the equilibrium. The only structural differences are mechanical translations of
the continuous-time objects: the infinitesimal generator $L$ becomes a one-period
transition (Markov) operator $\Tr$, the discount \emph{rate} $\rho$ becomes a
discount \emph{factor} $\beta$, the Kolmogorov forward equation
$\partial_t g = L^{*}[g]$ becomes the one-step push-forward $g_{t+1}=\Tstar[g_t]$,
and ``sums of generators'' ($L_1+L_2$) become ``products of transition operators''
($\Tr_1\Tr_2$). We verify that the discrete-time Master Equation collapses to
Bilal's continuous-time Master Equation as the period length $\Delta t\to 0$, and
we derive the discrete-time FAME, including aggregate shocks.
\end{abstract}

\tableofcontents

\section{Setup}

Time is discrete, $t=0,1,2,\dots$. There is a unit mass of households. A household
is described by the idiosyncratic state $x=(a,y)$, where $a\ge 0$ are asset holdings
(subject to a borrowing limit $a\ge\underline a$, normalized to $\underline a=0$)
and $y$ is idiosyncratic labor productivity. Let $\X=[\,0,\infty)\times Y$ denote
the idiosyncratic state space.

\paragraph{Idiosyncratic productivity.}
Productivity $y_t$ is an exogenous, stationary Markov process, independent across
households, with transition kernel $P(y,\dd y')=\Pr(y_{t+1}\in\dd y'\mid y_t=y)$.
Define the associated \emph{income transition operator} (a Markov, i.e.\ expectation,
operator)
\[
  (P\varphi)(y)\;\equiv\;\E\!\left[\varphi(y_{t+1})\mid y_t=y\right]
  \;=\;\int \varphi(y')\,P(y,\dd y').
\]
In continuous time Bilal uses the \emph{generator} $M(y)$ of the productivity
process (e.g.\ $M(y)[V]=\mu(y)V'(y)+\tfrac{\sigma(y)^2}{2}V''(y)$ for a diffusion).
Its discrete-time counterpart is the one-period operator $P$, related to a generator
by $P=\mathrm{Id}+\Delta t\,M+o(\Delta t)$.

\paragraph{Preferences and budget.}
Households have time-separable preferences with period utility $u$ (e.g.\ CRRA,
$u(c)=\tfrac{c^{1-\gamma}-1}{1-\gamma}$) and discount factor $\beta\in(0,1)$. They
save in a single risk-free asset and face the flow budget constraint
\begin{equation}\label{eq:budget}
  c_t + a_{t+1} \;=\; (1+r_t)\,a_t + w_t\,y_t,
  \qquad a_{t+1}\ge 0,
\end{equation}
where $r_t$ is the net interest rate and $w_t$ the wage. (The discount factor maps
to the continuous-time rate by $\beta=e^{-\rho\Delta t}$; the gross return
$1+r_t$ maps to $1+r_t\Delta t$.)

\paragraph{Technology and prices.}
A representative firm operates $Y_t=Z\,K_t^{\alpha}N_t^{1-\alpha}$, renting capital
$K_t$ (equal to aggregate assets) and hiring labor $N_t$ (equal to aggregate
efficiency units). With no depreciation, profit maximization and market clearing
pin prices to the current distribution $g_t$ of $x=(a,y)$:
\begin{equation}\label{eq:prices}
  r_t = R(g_t)\equiv \alpha Z\!\left(\frac{K(g_t)}{N(g_t)}\right)^{\alpha-1},
  \qquad
  w_t = W(g_t)\equiv (1-\alpha)Z\!\left(\frac{K(g_t)}{N(g_t)}\right)^{\alpha},
\end{equation}
where $K(g)=\int a\,g(x)\dd x$ and $N(g)=\int y\,g(x)\dd x$. As in Bilal's
equation~(2), $R$ and $W$ are functionals of the distribution.

\section{Sequential equilibrium}

\paragraph{Household problem.}
Taking the price sequences $\{r_t,w_t\}$ as given, the household solves
\begin{equation}\label{eq:seq-bellman}
  V_t(a,y) \;=\; \max_{a'\ge 0}\;
  u\!\big((1+r_t)a + w_t y - a'\big)
  \;+\; \beta\,\E\!\left[V_{t+1}(a',y')\,\middle|\,y\right].
\end{equation}
Let $\hat a_{t}(x)\equiv \hat s_t(x)$ denote the optimal saving policy and
$\hat c_t(x)=(1+r_t)a+w_ty-\hat s_t(x)$ optimal consumption. The interior
optimality (Euler) condition is
\begin{equation}\label{eq:euler}
  u'(\hat c_t(x))\;=\;\beta\,(1+r_{t+1})\,\E\!\left[u'(\hat c_{t+1}(a',y'))\mid y\right],
\end{equation}
with complementary slackness at the borrowing constraint $a'=0$.

\paragraph{Controlled transition operator.}
Bundle the deterministic asset transition $a'=\hat s_t(x)$ with the stochastic
income transition $P$ into the one-period \emph{controlled transition operator}
$\Tr_t$ acting on functions of $x$,
\begin{equation}\label{eq:Top}
  (\Tr_t\varphi)(x)\;\equiv\;\E\!\left[\varphi(x_{t+1})\mid x_t=x\right]
  \;=\;\int \varphi\big(\hat s_t(x),\,y'\big)\,P(y,\dd y'),
\end{equation}
the discrete-time analogue of the generator $L_t(x,c)$ in Bilal's equation~(1).
Its adjoint $\Tstar_t$ pushes a distribution one period forward.

\paragraph{Law of motion of the distribution.}
The continuous-time Kolmogorov forward equation
$\partial_t g_t=L_t^{*}[g_t]$ (Bilal eq.~3) becomes, in discrete time, the
one-step \emph{push-forward}
\begin{equation}\label{eq:kfe}
  g_{t+1}(x') \;=\; \big(\Tstar_t\,g_t\big)(x')
  \;=\; \int \mathbf 1\{\hat s_t(a,y)=a'\}\,P(y,\dd y')\,g_t(a,y)\dd a .
\end{equation}
Equation~\eqref{eq:kfe} is the discrete-time Kolmogorov forward equation: it simply
transports mass according to the saving policy and the income chain.

\paragraph{Equilibrium.} Given $g_0$, a sequential equilibrium is a path
$\{V_t,\hat s_t,r_t,w_t,g_t\}$ such that \eqref{eq:seq-bellman} holds with prices
\eqref{eq:prices} and the distribution evolves by \eqref{eq:kfe}. As in continuous
time, the difficulty is the forward--backward structure: \eqref{eq:seq-bellman} is
solved backward, \eqref{eq:kfe} forward, and prices are the fixed point that clears
markets each period.

\section{The discrete-time Master Equation}

Following Bilal, make the distribution an explicit state variable. Substitute prices
$r_t=R(g_t)$, $w_t=W(g_t)$ into the Bellman equation and replace the time index by
an explicit dependence on the distribution, $V_t(x)\equiv V(x,g_t)$. Knowing the law
of motion \eqref{eq:kfe} means the household can forecast next period's distribution:
write the one-step distribution map
\begin{equation}\label{eq:Phi}
  \Phi(g)\;\equiv\;\Tstar(\,\cdot\,,g)[\,g\,],
\end{equation}
i.e.\ $\Phi(g)$ is the distribution obtained by applying \eqref{eq:kfe} once, using
the policy that is optimal given the current distribution $g$. Then the household
problem becomes fully recursive:
\begin{empheq}[box=\fbox]{equation}\label{eq:master}
  V(x,g) \;=\; \max_{a'\ge 0}\;
  u\!\big((1+R(g))a + W(g)y - a'\big)
  \;+\; \beta\,\int V\!\big((a',y'),\,\Phi(g)\big)\,P(y,\dd y').
\end{empheq}
Equation~\eqref{eq:master} is the \textbf{discrete-time Master Equation}. It is the
exact discrete-time counterpart of Bilal's continuous-time Master Equation
\[
  \rho V(x,g)=\max_{c} u(c) + L(x,c,g)[V]
  + \int \frac{\partial V}{\partial g}(x,x',g)\,L^{*}(x',g)[g]\dd x'.
  \tag{Bilal, eq.~7}
\]
The two continuation terms of the continuous-time equation---the partial-equilibrium
generator $L[V]$ (evolution of \emph{own} state) and the term
$\int \partial_g V\,L^{*}[g]$ (evolution of the \emph{distribution})---are jointly
encoded in the single discrete-time continuation
$\beta\!\int V((a',y'),\Phi(g))P(y,\dd y')$: the conditional expectation over
$(a',y')$ captures the evolution of the household's own state, while the argument
$\Phi(g)$ captures the evolution of the aggregate distribution. As in continuous
time, \eqref{eq:master} is the \emph{only} equation that needs to be solved: the
distributional fixed point has been absorbed into the value function, so the
representation is block-recursive.

\subsection{Continuous-time limit (consistency check)}

Let the period length be $\Delta t$, set $\beta=e^{-\rho\Delta t}$, scale the budget
as a flow $a_{t+\Delta t}=a_t+(r_ta_t+w_ty_t-c_t)\Delta t$, and expand the operators
to first order in $\Delta t$:
\[
  \Tr = \mathrm{Id}+\Delta t\,L + o(\Delta t),
  \qquad
  \Phi(g)=g+\Delta t\,L^{*}(g)[g]+o(\Delta t).
\]
The Master Equation \eqref{eq:master} (with flow utility $u(c)\,\Delta t$) becomes
\[
  V(x,g)=\max_c u(c)\,\Delta t
  + (1-\rho\Delta t)\,\big(\mathrm{Id}+\Delta t\,L_x\big)
  \Big[\,V(x,g)+\Delta t\!\int \tfrac{\partial V}{\partial g}(x,\xi,g)\,L^{*}(\xi,g)[g]\dd\xi\Big].
\]
Cancelling $V(x,g)$, dividing by $\Delta t$ and letting $\Delta t\to 0$ recovers
exactly Bilal's equation~(7). Thus the discrete-time formulation is consistent with,
and reduces to, the continuous-time Master Equation---confirming that ``the same
methodology naturally works in discrete time.''

\section{The discrete-time FAME}

We now perturb \eqref{eq:master} around a deterministic steady state, following
Bilal's First-order Approximation to the Master Equation (FAME).

\paragraph{Steady state.}
A steady state is $(\Vss,\gss)$ with prices $r^{\mathrm{ss}}=R(\gss)$,
$w^{\mathrm{ss}}=W(\gss)$, policy $\sss(x)$, consumption $\css(x)$, solving
\[
  \Vss(x)=\max_{a'\ge0} u\big((1+r^{\mathrm{ss}})a+w^{\mathrm{ss}}y-a'\big)+\beta\,(\Tr\,\Vss)(x),
  \qquad \gss=\Tstar\gss,
\]
where $\Tr$, $\Tstar$ are the steady-state transition operator \eqref{eq:Top} and its
adjoint, evaluated at $\sss$.

\paragraph{Impulse Value ansatz.}
Let $h\equiv g-\gss$ be a small distributional impulse. To first order the value
function is linear in $h$:
\begin{equation}\label{eq:ansatz}
  V(x,\gss+h)\;=\;\Vss(x)\;+\;\int v(x,\xi)\,h(\xi)\dd\xi \;+\;o(\|h\|),
\end{equation}
where $v(x,\xi)=\dfrac{\partial V}{\partial g}(x,\xi,\gss)$ is the \emph{Impulse
Value}: the general-equilibrium effect on the value of a household at $x=(a,y)$ of
adding an infinitesimal mass of households at $\xi=(a_\xi,y_\xi)$. This is identical
to the continuous-time object; only the equation it satisfies differs.

\paragraph{Derivation.}
Substitute \eqref{eq:ansatz} into \eqref{eq:master} and linearize in $h$, using the
envelope theorem (the saving policy is held at its optimum, so its first-order
variation drops out). Three first-order channels appear.

\emph{(i) Direct price impact.} Adding mass at $\xi$ changes $K$ by $a_\xi$ and $N$ by
$y_\xi$, hence moves prices by $\dd r=R_K a_\xi+R_N y_\xi$ and
$\dd w=W_K a_\xi+W_N y_\xi$ (derivatives of \eqref{eq:prices} at $\gss$). The induced
change in disposable income, valued at the steady-state marginal utility, is
$u'(\css(x))\,D(x,\xi)$ with the \emph{price-impact function}
\begin{equation}\label{eq:D}
  D(x,\xi)\;=\;\big(R_K a+W_K y\big)\,a_\xi+\big(R_N a+W_N y\big)\,y_\xi,
\end{equation}
exactly Bilal's $D(x,x')=(R_0 a_\xi+R_1 y_\xi)a+(W_0 a_\xi+W_1 y_\xi)y$ with
$R_0=R_K,\,R_1=R_N,\,W_0=W_K,\,W_1=W_N$.

\emph{(ii) Continuation through transitions.} The continuation
$\beta\int v(x^{+},\xi)\,\dd[\Phi]$ expands using $\Phi(\gss+h)=\gss+\Tstar h+\Gop h$,
where $\Tstar h$ is the partial-equilibrium push-forward of the impulse (policies
fixed) and $\Gop$ is the general-equilibrium response of the push-forward to
policy changes (defined below). The push-forward part, after taking adjoints,
transports the \emph{second} argument of $v$ one period forward; the expectation
over the household's own next state transports the \emph{first} argument.

\emph{(iii) General-equilibrium savings response.} The price change makes every
household $\zeta$ revise its saving by $\delta s(\zeta,\xi)$, which reshapes next
period's distribution through the kernel $G(\cdot,\xi)$ of $\Gop$.

Collecting terms, the \textbf{discrete-time deterministic FAME} is
\begin{empheq}[box=\fbox]{equation}\label{eq:fame}
  v(x,\xi)\;=\;\underbrace{u'(\css(x))\,D(x,\xi)}_{\text{direct price impact}}
  \;+\;\beta\,\underbrace{\big(\Tr_1\Tr_2\,v\big)(x,\xi)}_{\substack{\text{continuation:}\\ \text{own state }x\text{ and impulse }\xi}}
  \;+\;\beta\,\underbrace{\Big(\Tr_1\!\big[v\star G\big]\Big)(x,\xi)}_{\substack{\text{GE: savings responses}\\ \text{of other households}}},
\end{empheq}
where $\Tr_1$ ($\Tr_2$) applies the transition operator \eqref{eq:Top} to the first
(second) argument of $v$, and
$(v\star G)(x,\xi)\equiv\int v(x,\zeta)\,G(\zeta,\xi)\dd\zeta$.

\paragraph{Comparison with continuous time.}
Bilal's continuous-time FAME (his eq.~9) reads
\[
  \rho v(x,x') = u'(\css(x))D(x,x')
   + L(x)[v(\cdot,x')] + L(x')[v(x,\cdot)] + \!\int\! v(x,x'')G(x'',x',v)\dd x''.
\]
The dictionary is exact:
\begin{itemize}[itemsep=2pt]
\item the discount factor $\beta$ replaces $1/\rho$ (more precisely $1-\rho\,\Delta t$);
\item the \emph{product} of transition operators $\Tr_1\Tr_2$ replaces the \emph{sum}
of generators $L(x)+L(x')$ (because
$\Tr_1\Tr_2=(\mathrm{Id}+\Delta t L_1)(\mathrm{Id}+\Delta t L_2)=\mathrm{Id}+\Delta t(L_1+L_2)+o(\Delta t)$);
\item the GE kernel $G$ is the discrete-time analogue of $G(x'',x',v)$ and likewise
depends on $v$ through the consumption (dMPC) response, so \eqref{eq:fame} is
\emph{quadratic} in $v$, just like the continuous-time FAME.
\end{itemize}

\paragraph{The general-equilibrium kernel.}
The per-period saving response of household $\zeta$ to a unit impulse at $\xi$ is the
change in resources minus the consumption response,
\[
  \delta s(\zeta,\xi)\;=\;D(\zeta,\xi)\;-\;\delta c(\zeta,\xi),
\]
where $\delta c$ is obtained by linearizing the Euler equation \eqref{eq:euler}
(a ``distributional MPC''). The kernel $G$ is then the linearized push-forward of the
displaced mass: in continuous time
$G(\zeta,\xi,v)=-\partial_{a_\zeta}\!\big[\gss(\zeta)\,\delta s(\zeta,\xi)\big]$ with
$\delta c=\tfrac{1}{u''(\css)}\,\partial_a v$; in discrete time the divergence
$-\partial_{a}(\cdot)$ is replaced by the linearization of the push-forward
\eqref{eq:kfe} with respect to the saving policy,
\[
  \big(\Gop h\big)(x')\;=\;\frac{\partial}{\partial g}\Big[\Tstar(\cdot,g)[g]\Big]_{\gss}\!\!\cdot h,
  \qquad
  \big(\Gop h\big)(x')=\int G(x',\xi)\,h(\xi)\dd\xi .
\]

\subsection{Matrix (Sylvester) form}

Discretize $x$ on a grid $\{x_i\}$, write $\mathbf V=[v(x_i,x_j)]$, let $T$ be the
steady-state transition matrix (row-stochastic), $T^{\top}$ its transpose,
$\widetilde D=[\,u'(\css_i)D(x_i,x_j)\,]$, and $G=G(\mathbf V)$ the discretized GE
kernel. Because $\Tr_1$ left-multiplies, $\Tr_2$ right-multiplies by $T^{\top}$, and
$v\star G$ is right-multiplication by $G$, the FAME \eqref{eq:fame} becomes the
\textbf{nonlinear (Riccati-type) Sylvester equation}
\begin{empheq}[box=\fbox]{equation}\label{eq:sylvester}
  \mathbf V \;=\; \widetilde D \;+\; \beta\,T\,\mathbf V\,\big(T^{\top}+G(\mathbf V)\big).
\end{empheq}
This is the discrete-time analogue of Bilal's continuous-time Sylvester equation
$\rho v = D + Lv + vL^{\top} + vGv$ and is solved by the same fast iterative schemes.

\subsection{Linearized law of motion of the distribution}

Perturbing \eqref{eq:kfe} once $v$ (hence $G$) is known gives the discrete-time
analogue of Bilal's eq.~(10),
\begin{equation}\label{eq:lom}
  h_{t+1}\;=\;\Tstar h_t \;+\; \Gop h_t
  \qquad\Longleftrightarrow\qquad
  h_{t+1}=\big(T^{\top}+G\big)\,h_t .
\end{equation}
The first term is the partial-equilibrium transport of the impulse; the second is the
general-equilibrium feedback through savings. Iterating \eqref{eq:lom} forward
delivers impulse responses without solving any further fixed point (block
recursivity). \textbf{Stability} is the discrete-time mirror of the continuous-time
condition: convergence back to steady state requires the spectral radius of
$T^{\top}+G$ to be strictly below $1$ (rather than the eigenvalues of $L^{*}+\Gop$
having strictly negative real part).

\section{The Huggett variant: an implicitly priced asset}\label{sec:huggett}

In Krusell--Smith the asset is productive capital and its price is the marginal
product, a \emph{technological} functional of the distribution. Huggett (1993) is the
polar case: a pure-exchange economy in which the asset is in fixed net supply and its
price is determined \emph{behaviorally}, by market clearing. The Master Equation
accommodates this with no change in logic---only the price functional is replaced.

\subsection{Setup}

There is no production: $y$ is an exogenous \emph{endowment} (still Markov with
operator $P$), so there is no wage. Households trade a one-period bond. Let $q_t$ be
the \emph{current bond price}: a unit of the bond bought at $t$ costs $q_t$ and pays
one unit of the good at $t+1$, so
\begin{equation}\label{eq:hug-q}
  \frac{1}{q_t}\;=\;1+r_t .
\end{equation}
Let $a_t$ denote the (matured) bond holdings carried into $t$, each paying one unit,
and $a_{t+1}$ the bonds purchased at $t$. The budget constraint is
\begin{equation}\label{eq:hug-budget}
  c_t + q_t\,a_{t+1} \;=\; a_t + y_t, \qquad a_{t+1}\ge \underline a .
\end{equation}
Writing the budget with $q_t$ rather than $r_t$ is the economically correct timing:
the period-$t$ saving choice trades at the \emph{current} price $q_t$, whereas the
realized return $r_t=1/q_t-1$ in \eqref{eq:hug-q} is the payoff on \emph{already-held}
bonds $a_t$ and is predetermined. Thus $r_t$ does not affect time-$t$ behavior; $q_t$
does, and $q_t$ is the price that clears the market. The bond is in fixed net supply
$B$ (pure Huggett sets $B=0$; $B>0$ represents government debt).

\paragraph{The asset-demand function.} To determine the price, treat the current bond
price $q$ as a \emph{free variable} rather than imposing $q=Q(g)$. For a candidate $q$,
let
\begin{equation}\label{eq:hug-demand}
  \mathcal A(x;\,q,\,g)\;\equiv\;\arg\max_{a'\ge\underline a}\;
  u\big(a+y-q\,a'\big)+\beta\!\int V\big((a',y'),\Phi(g)\big)\,P(y,\dd y')
\end{equation}
denote the household's bond demand when the \emph{current} trade is at price $q$ and
the continuation is taken at the equilibrium law of motion $\Phi(g)$. This is the
\emph{asset-demand function}: because $q$ enters the current budget, $\mathcal A$
genuinely depends on it, with $\partial_q\mathcal A<0$ in the usual case (demand falls
as the bond gets more expensive). The \emph{equilibrium} policy is the special case in
which the trade occurs at the market-clearing price,
\[
  \hat a(x,g)\;=\;\mathcal A\big(x;\,Q(g),\,g\big),
\]
which is why the equilibrium policy $\hat a(x,g)$ alone carries no free price: it has
$q=Q(g)$ already substituted. The bond price $Q(g)$ is the value that clears the
market,
\begin{equation}\label{eq:hug-clearing}
  \int \mathcal A(x;\,q,\,g)\,g(x)\dd x \;=\; B
  \qquad\Longrightarrow\qquad
  q \;=\; Q(g)\ \ \text{(defined \emph{implicitly}).}
\end{equation}
Since $\partial_q\mathcal A\neq 0$, \eqref{eq:hug-clearing} is a genuine
(one-dimensional) equation in $q$.

\subsection{Master Equation}

Nothing changes in the recursive logic: $q=Q(g)$ is a (Fr\'echet-differentiable)
functional of the distribution, so Bilal's Assumption~2 is satisfied with an
implicitly defined functional. The discrete-time Master Equation \eqref{eq:master}
becomes
\begin{equation}\label{eq:hug-master}
  V(x,g) \;=\; \max_{a'\ge \underline a}\;
  u\!\big(a + y - Q(g)\,a'\big)
  \;+\; \beta\,\int V\!\big((a',y'),\,\Phi(g)\big)\,P(y,\dd y'),
\end{equation}
with the same one-step distribution map $\Phi(g)=\Tstar(\cdot,g)[g]$. Steady state is
a pair $(\Vss,\gss)$ with $q^{\mathrm{ss}}=Q(\gss)$ clearing \eqref{eq:hug-clearing}.

\subsection{FAME}

The Impulse Value ansatz \eqref{eq:ansatz} and the FAME \eqref{eq:fame} are unchanged
in form. Two things are different relative to Krusell--Smith.

\paragraph{(i) A single bond-price channel.} Income is an exogenous endowment, so there
is no wage. The only price is $q=Q(g)$, and it enters the household's value only
through the cost $q\,a'$ of new bonds. By the envelope theorem, the direct effect of a
price change $\dd q=Q'(\gss)[\delta_\xi]$ on the value, valued at steady-state marginal
utility, gives the price-impact function
\begin{equation}\label{eq:hug-D}
  D(x,\xi) \;=\; -\,\hat a^{\mathrm{ss}}(x)\;Q'(\gss)[\delta_\xi],
\end{equation}
i.e.\ a household that is buying bonds ($\hat a^{\mathrm{ss}}(x)>0$) is hurt by a higher
bond price, in proportion to its bond purchases. (Contrast Krusell--Smith, where the
price impact was proportional to asset \emph{holdings} $a$ and worked through both $r$
and $w$.) Here $Q'(\gss)[\delta_\xi]$ is the marginal effect on the equilibrium bond
price of placing an infinitesimal unit of mass at $\xi=(a_\xi,y_\xi)$.

\paragraph{(ii) The price impact is implicit and endogenous to $v$.}
Apply the implicit function theorem to the market-clearing condition
\eqref{eq:hug-clearing}, differentiating in the direction $h$. The distribution
$g$ enters \eqref{eq:hug-clearing} both directly (as the measure $\int\mathcal A\,g$)
and through the continuation $\Phi(g)$ inside $\mathcal A$, so
\begin{equation}\label{eq:hug-Qprime}
  Q'(\gss)[h]
  \;=\;-\,\frac{\displaystyle \int \hat a^{\mathrm{ss}}(x)\,h(x)\dd x
        \;+\;\int \frac{\partial \mathcal A}{\partial g}\!\cdot h\;\gss(x)\dd x}
       {\displaystyle \int \frac{\partial \mathcal A}{\partial q}(x)\;\gss(x)\dd x },
  \qquad \hat a^{\mathrm{ss}}(x)=\mathcal A(x;q^{\mathrm{ss}},\gss).
\end{equation}
The denominator $\int \partial_q\mathcal A\;\gss<0$ is the \emph{slope of the aggregate
bond-demand curve}; the price moves \emph{more} the \emph{less} elastic aggregate
demand is. In the numerator, the first term is the mechanical composition effect
(adding mass at $\xi$ contributes $\hat a^{\mathrm{ss}}(\xi)$ to demand directly), and
the second, $\int \partial_g\mathcal A\cdot h\,\gss$, is the response of demand to the
change in the \emph{future} path of prices that $h$ induces through $\Phi(g)$---the
channel that carries the Impulse Value $v$ into the price. In Krusell--Smith the
denominator was instead the closed-form curvature of the marginal product, so
$D(x,\xi)$ was known once $\gss$ was known, \emph{independently} of $v$. In Huggett,
$\partial_g\mathcal A$ depends on the household's valuation of distributional
change---i.e.\ on $v$ itself, through the dMPC. Hence $D=D(v)$, and the FAME becomes the
fixed point
\begin{equation}\label{eq:hug-sylvester}
  \mathbf V \;=\; \widetilde D(\mathbf V) \;+\; \beta\,T\,\mathbf V\,\big(T^{\top}+G(\mathbf V)\big),
\end{equation}
where now \emph{both} $\widetilde D$ (via the market-clearing elasticity
\eqref{eq:hug-Qprime}) and $G$ depend on $\mathbf V$. The GE kernel $G$ retains the
form of Section~4, with the saving response $\delta s(\zeta,\xi)=D(\zeta,\xi)-\delta
c(\zeta,\xi)$ now driven by the single implicit price \eqref{eq:hug-Qprime}.
Everything else---block recursivity, the linearized law of motion
\eqref{eq:lom}, and the stability criterion (spectral radius of $T^{\top}+G$ below
$1$)---carries over unchanged.

\paragraph{Summary of the contrast.} Huggett is nested in the same Master Equation as
Krusell--Smith: one passes from one to the other simply by replacing the
marginal-product price functional with the market-clearing functional
\eqref{eq:hug-clearing}. The substantive consequence is that the price impact
$D$ is no longer a closed-form steady-state object but is itself part of the FAME fixed
point, because the market-clearing price responds to the same household behavior that
the Impulse Value encodes.

\section{Aggregate shocks}

Let aggregate productivity be $Z_t=\bar Z\,e^{\varepsilon z_t}$, with $z_t$ a
stationary AR(1),
\[
  z_{t+1}=\phi\,z_t+\sigma\,\eta_{t+1},\qquad |\phi|<1,\ \ \eta_{t+1}\sim(0,1),
\]
and define the aggregate expectation operator
$(\Aop\psi)(z)\equiv\E[\psi(z_{t+1})\mid z_t=z]$, the discrete-time analogue of the
generator $A(z)$ in Bilal's eq.~(11). Prices \eqref{eq:prices} now depend on $z$
through $Z$. The \textbf{Master Equation with aggregate shocks} is
\begin{equation}\label{eq:master-agg}
  V(x,z,g)=\max_{a'\ge0}\;
  u(c)
  +\beta\,\E\Big[\,V\big((a',y'),\,z',\,\Phi(g,z)\big)\ \big|\ y,z\,\Big],
\end{equation}
where $\Phi(g,z)$ is the one-step distribution map evaluated at current aggregate
state $z$ and the expectation is over $(y',z')$.

\paragraph{First-order expansion.}
With both the distributional deviation and the shock scale of order $\varepsilon$,
\[
  V(x,z,\gss+\varepsilon h)=\Vss(x)+\varepsilon\Big\{\int v(x,\xi)h(\xi)\dd\xi+\omega(x,z)\Big\}+o(\varepsilon),
\]
where the deterministic Impulse Value $v$ still solves the deterministic FAME
\eqref{eq:fame} (certainty equivalence / block recursivity), and $\omega(x,z)$ is the
\emph{stochastic Impulse Value}. Substituting into \eqref{eq:master-agg} yields the
\textbf{discrete-time stochastic FAME}
\begin{empheq}[box=\fbox]{equation}\label{eq:sfame}
  \omega(x,z)=\underbrace{z\,\Omega(x)\,u'(\css(x))}_{\text{direct price impact of }z}
  +\beta\,\underbrace{\big(\Aop\,\Tr_1\,\omega\big)(x,z)}_{\substack{\text{continuation:}\\ \text{own state and }z}}
  +\beta\,\underbrace{\Tr_1\!\Big[\,v\star S(\cdot,z)\,\Big](x)}_{\substack{\text{GE: savings responses}\\ \text{to the aggregate shock}}},
\end{empheq}
where $\Omega(x)=R_Z\,a+W_Z\,y$ collects the direct effect of $Z$ on prices
(the analogue of Bilal's $\Omega(x)=R_2a+W_2y$), and $S(\cdot,z)$ is the GE savings
response to the aggregate shock. Given $v$, equation \eqref{eq:sfame} is \emph{linear}
in $\omega$ (the deterministic FAME was quadratic in $v$), exactly as in continuous
time, and reduces to a linear system after discretization.

\paragraph{Law of motion with aggregate shocks.}
The linearized distribution evolves as the discrete-time analogue of Bilal's eq.~(15),
\begin{equation}\label{eq:lom-agg}
  h_{t+1}=\big(T^{\top}+G\big)\,h_t + S\,z_t,
\end{equation}
so that iterating \eqref{eq:lom-agg} with $\{z_t\}$ produces aggregate impulse
responses and characterizes the stochastic steady state.

\section{Summary: the continuous-to-discrete dictionary}

\begin{center}
\renewcommand{\arraystretch}{1.35}
\begin{tabular}{l l l}
\hline
\textbf{Object} & \textbf{Continuous time (Bilal)} & \textbf{Discrete time (this note)}\\
\hline
Discounting & rate $\rho$ & factor $\beta=e^{-\rho\Delta t}$\\
Own-state evolution & generator $L(x,c,g)$ & transition operator $\Tr$, eq.~\eqref{eq:Top}\\
Income process & generator $M(y)$ & Markov operator $P$\\
Distribution law of motion & $\partial_t g=L^{*}[g]$ & $g_{t+1}=\Tstar[g_t]=\Phi(g_t)$\\
Master Equation & eq.~(7) & eq.~\eqref{eq:master}\\
Two continuation terms $L_1+L_2$ & sum of generators & product $\Tr_1\Tr_2$\\
FAME & $\rho v=D+Lv+vL^{\top}+vGv$ & $\mathbf V=\widetilde D+\beta T\mathbf V(T^{\top}+G)$\\
LoM of impulse & $\partial_t h=L^{*}h+\Gop h$ & $h_{t+1}=(T^{\top}+G)h_t$\\
Stability & $\mathrm{Re}\,\mathrm{eig}<0$ & spectral radius $<1$\\
Stochastic FAME & eq.~(14), linear in $\omega$ & eq.~\eqref{eq:sfame}, linear in $\omega$\\
\hline
\end{tabular}
\end{center}

\noindent
The discrete-time FAME is precisely the analytic, ``linearize-then-discretize''
foundation for state-space methods in discrete time (Reiter, 2009; Ahn et al., 2018)
and is the natural companion of the discrete-time perturbation techniques of
Bhandari et al.\ (2023) and the sequence-space Jacobian approach of Auclert et al.\
(2021).

\end{document}
